\documentclass[../DoAn.tex]{subfiles}
\begin{document}

\section{Giới thiệu chung về kiến trúc và công nghệ}

Việc lựa chọn công nghệ và kiến trúc hệ thống là một trong những quyết định nền tảng, có ảnh hưởng sâu sắc đến sự thành công, hiệu quả vận hành và khả năng phát triển trong tương lai của một dự án phần mềm. Đối với một ứng dụng học TOEIC trực tuyến, nơi tính tương tác, độ tin cậy và trải nghiệm người dùng được đặt lên hàng đầu, việc xây dựng một hệ thống vững chắc là yếu tố sống còn.

Dự án này được thiết kế và phát triển dựa trên kiến trúc \texttt{client-server} hiện đại. Theo đó, hai thành phần chính của hệ thống được tách biệt rõ ràng:
\begin{itemize}
    \item \textbf{Frontend (Client-side):} Được xây dựng dưới dạng một Single Page Application (\texttt{SPA}), sử dụng thư viện \textbf{React} kết hợp với \textbf{TypeScript} để đảm bảo mã nguồn có cấu trúc chặt chẽ và an toàn. Giao diện người dùng được thiết kế với \textbf{Tailwind CSS} và thư viện component \textbf{Shadcn/UI}, mang lại sự linh hoạt, nhất quán và khả năng tùy biến cao.
    \item \textbf{Backend (Server-side):} Được phát triển bằng \textbf{Spring Boot}, một framework mạnh mẽ của hệ sinh thái Java, theo kiến trúc \textbf{RESTful API}. Hệ thống sử dụng \textbf{MySQL} làm hệ quản trị cơ sở dữ liệu quan hệ để lưu trữ dữ liệu một cách bền vững và \textbf{Cloudinary} để quản lý file media trên cloud. Cơ chế bảo mật được triển khai dựa trên \textbf{JSON Web Token (JWT)}, một tiêu chuẩn ngành cho việc xác thực và phân quyền an toàn.
\end{itemize}

Sự lựa chọn này không phải là ngẫu nhiên mà đến từ việc phân tích kỹ lưỡng các tiêu chí quan trọng. Các công nghệ được chọn đều có cộng đồng phát triển đông đảo, nguồn tài liệu phong phú và được cập nhật liên tục, đảm bảo dự án không bị lạc hậu và dễ dàng tìm kiếm sự hỗ trợ khi cần. Hơn nữa, kiến trúc tách biệt giữa frontend và backend không chỉ giúp tối ưu hóa quy trình làm việc song song của các nhóm phát triển mà còn tạo ra một nền tảng linh hoạt, sẵn sàng cho việc mở rộng hệ thống trong tương lai, chẳng hạn như phát triển ứng dụng di động hoặc tích hợp với các dịch vụ của bên thứ ba.

\section{Frontend: Xây dựng giao diện người dùng hiện đại}

\subsection{React và TypeScript: Nền tảng cho giao diện linh hoạt và an toàn}

\textbf{React}, thư viện JavaScript do Facebook phát triển, đã trở thành một tiêu chuẩn de facto trong việc xây dựng giao diện người dùng (UI) cho các ứng dụng web hiện đại. Trái tim của React là mô hình component, cho phép các nhà phát triển chia nhỏ giao diện phức tạp thành các thành phần độc lập, có thể tái sử dụng và dễ dàng quản lý. Mỗi component có vòng đời (lifecycle) và trạng thái (state) riêng, giúp việc xử lý dữ liệu và logic giao diện trở nên mạch lạc. Điểm ưu việt nhất của React chính là việc sử dụng Virtual DOM (DOM ảo). Thay vì thao tác trực tiếp lên DOM thật của trình duyệt – một quá trình tốn kém về hiệu năng – React tạo ra một bản sao của DOM trong bộ nhớ. Khi trạng thái của một component thay đổi, React sẽ tính toán sự khác biệt (diffing) và chỉ cập nhật những phần cần thiết trên DOM thật, giúp tối ưu hóa đáng kể tốc độ render và mang lại trải nghiệm mượt mà cho người dùng.

Để nâng cao chất lượng và độ tin cậy của mã nguồn, dự án tích hợp \textbf{TypeScript}. TypeScript là một siêu tập hợp (superset) của JavaScript, được phát triển bởi Microsoft, bổ sung hệ thống kiểm tra kiểu tĩnh (static typing). Trong một dự án lớn như hệ thống học tập trực tuyến với nhiều logic phức tạp, việc sử dụng JavaScript thuần túy có thể dẫn đến các lỗi \texttt{runtime} khó lường do sai kiểu dữ liệu. TypeScript giải quyết triệt để vấn đề này bằng cách cho phép khai báo kiểu cho biến, tham số hàm và các cấu trúc dữ liệu. Trình biên dịch của TypeScript sẽ phát hiện các lỗi này ngay trong quá trình phát triển (\texttt{compile-time}), trước cả khi mã được chạy. Việc này không chỉ giúp giảm thiểu lỗi, tăng tính ổn định của ứng dụng mà còn làm cho mã nguồn trở nên tự miêu tả (self-documenting), dễ đọc và dễ bảo trì hơn, đặc biệt hữu ích khi làm việc nhóm.

\subsection{Tailwind CSS và Shadcn/UI: Tối ưu hóa quy trình thiết kế giao diện}

\textbf{Tailwind CSS} được lựa chọn làm framework CSS chính cho dự án, theo triết lý "utility-first". Khác với các framework truyền thống như Bootstrap vốn cung cấp các component dựng sẵn (như \texttt{.card}, \texttt{.button}), Tailwind cung cấp một bộ sưu tập lớn các lớp tiện ích cấp thấp (ví dụ: \texttt{flex}, \texttt{pt-4}, \texttt{text-center}). Cách tiếp cận này cho phép các nhà phát triển xây dựng các thiết kế hoàn toàn tùy chỉnh trực tiếp trong mã HTML/JSX mà không cần phải rời khỏi trình soạn thảo hay viết các file CSS riêng biệt. Điều này không chỉ tăng tốc độ phát triển mà còn đảm bảo tính nhất quán trong toàn bộ hệ thống. Hơn nữa, khi build sản phẩm, Tailwind sẽ tự động quét mã nguồn và loại bỏ tất cả các lớp CSS không được sử dụng, giúp dung lượng file CSS cuối cùng siêu nhỏ, góp phần cải thiện tốc độ tải trang.

Để đẩy nhanh hơn nữa quá trình xây dựng giao diện, dự án sử dụng \textbf{Shadcn/UI}. Đây không phải là một thư viện component thông thường, mà là một bộ sưu tập các component có thể tái sử dụng, được xây dựng trên nền tảng Tailwind CSS và Radix UI. Điểm đặc biệt của Shadcn/UI là bạn có thể sao chép và dán mã nguồn của các component trực tiếp vào dự án của mình, cho phép toàn quyền kiểm soát và tùy biến chúng. Điều này giúp đội ngũ phát triển vừa có được những component chất lượng cao, chuẩn mực về thiết kế, hỗ trợ responsive và accessibility, vừa không bị ràng buộc bởi các quyết định thiết kế của thư viện.

\subsection{Quản lý trạng thái và điều hướng: Đảm bảo luồng dữ liệu thông suốt}

Trong một ứng dụng \texttt{SPA} phức tạp, việc quản lý trạng thái (state management) là một thách thức lớn. Dự án sử dụng các công cụ mạnh mẽ như \textbf{Redux Toolkit} hoặc \textbf{React Context API} để giải quyết bài toán này. Redux Toolkit, phiên bản cải tiến của Redux, cung cấp một cách thức tập trung để lưu trữ toàn bộ trạng thái của ứng dụng trong một "store" duy nhất. Điều này giúp việc theo dõi và gỡ lỗi luồng dữ liệu trở nên dễ dàng hơn. Trong khi đó, React Context API là một giải pháp tích hợp sẵn của React, phù hợp cho việc chia sẻ các trạng thái đơn giản hơn giữa các component mà không cần thông qua \texttt{props drilling} (truyền props qua nhiều cấp).

Việc điều hướng trong ứng dụng được quản lý bởi \textbf{React Router}. Thư viện này cho phép tạo ra các tuyến đường (routes) ảo, ánh xạ các URL trên thanh địa chỉ của trình duyệt với các component tương ứng. Nhờ đó, người dùng có thể di chuyển giữa các trang (ví dụ: từ trang danh sách bài học sang trang chi tiết một bài thi) mà không cần tải lại toàn bộ trang web. Điều này không chỉ tạo ra trải nghiệm liền mạch, tức thời mà còn giúp duy trì trạng thái của ứng dụng, chẳng hạn như dữ liệu đang được nhập trong một biểu mẫu.

\section{Backend: Trái tim của hệ thống}

\subsection{Spring Boot và kiến trúc RESTful API: Xây dựng dịch vụ backend mạnh mẽ}

Phần backend của hệ thống được xây dựng trên nền tảng \textbf{Spring Boot}, một phần mở rộng của Spring Framework. Spring Boot được ưa chuộng trong cộng đồng Java vì khả năng đơn giản hóa quá trình phát triển ứng dụng. Nó tuân thủ nguyên tắc "convention over configuration" (ưu tiên quy ước hơn cấu hình), cung cấp cơ chế \texttt{auto-configuration} thông minh, giúp giảm thiểu đáng kể lượng mã boilerplate và các tệp cấu hình XML phức tạp. Điều này cho phép lập trình viên tập trung vào việc xây dựng logic nghiệp vụ thay vì sa đà vào các thiết lập ban đầu.

Hệ thống được thiết kế theo kiến trúc \textbf{RESTful API}, một tiêu chuẩn kiến trúc để thiết kế các ứng dụng mạng. REST (Representational State Transfer) sử dụng các phương thức HTTP tiêu chuẩn (\texttt{GET}, \texttt{POST}, \texttt{PUT}, \texttt{DELETE}) để biểu diễn các hành động trên tài nguyên. Ví dụ, một yêu cầu \texttt{GET /api/users/123} sẽ trả về thông tin của người dùng có ID 123. Kiến trúc này giúp tách biệt hoàn toàn logic của backend khỏi giao diện của frontend. Sự tách biệt này mang lại nhiều lợi ích: khả năng bảo trì tốt hơn, cho phép phát triển độc lập giữa hai đội ngũ, và quan trọng nhất là khả năng mở rộng hệ thống để hỗ trợ nhiều loại client khác nhau (web app, mobile app, desktop app) trong tương lai chỉ bằng cách gọi cùng một bộ API.

\subsection{Xác thực và bảo mật bằng JWT: Bảo vệ dữ liệu người dùng}

Để đảm bảo an toàn cho hệ thống và bảo vệ dữ liệu nhạy cảm của người dùng, dự án áp dụng cơ chế xác thực dựa trên \textbf{JSON Web Token (JWT)}. Đây là một phương pháp xác thực không trạng thái (stateless), hiện đại và an toàn.

Quy trình hoạt động chi tiết như sau:
\begin{itemize}
    \item \textbf{Đăng nhập và tạo Token:} Khi người dùng cung cấp thông tin đăng nhập (username/password) chính xác, backend sẽ xác thực và tạo ra một chuỗi JWT. Token này bao gồm ba phần: Header (chứa thuật toán mã hóa), Payload (chứa thông tin người dùng như ID, vai trò) và Signature (chữ ký số để xác minh tính toàn vẹn).
    \item \textbf{Gửi và lưu trữ Token:} Token này được gửi về cho client. Frontend sẽ lưu trữ token này một cách an toàn, thường là trong \texttt{localStorage}, \texttt{sessionStorage} hoặc \texttt{cookie} (với cờ \texttt{HttpOnly} và \texttt{Secure} để tăng cường bảo mật).
    \item \textbf{Xác thực các yêu cầu sau:} Với mỗi yêu cầu cần xác thực gửi lên server sau đó (ví dụ: lấy thông tin cá nhân, nộp bài thi), client sẽ đính kèm JWT trong header \texttt{Authorization} theo chuẩn \texttt{Bearer <token>}.
    \item \textbf{Xác minh Token phía Backend:} Backend, thông qua một bộ lọc (filter) bảo mật, sẽ chặn mọi yêu cầu, giải mã JWT và kiểm tra chữ ký của nó bằng một khóa bí mật (\texttt{secret key}) chỉ server biết. Nếu token hợp lệ và chưa hết hạn, yêu cầu sẽ được xử lý. Ngược lại, server sẽ trả về lỗi \texttt{401 Unauthorized}.
\end{itemize}

Phương pháp này giúp giảm tải cho server vì không cần phải lưu trữ thông tin phiên của người dùng, làm cho hệ thống dễ dàng mở rộng theo chiều ngang (horizontal scaling).

\subsection{Cơ sở dữ liệu MySQL và ORM: Quản lý dữ liệu hiệu quả}

\textbf{MySQL} được chọn làm hệ quản trị cơ sở dữ liệu quan hệ (RDBMS) cho dự án. MySQL nổi tiếng về sự ổn định, hiệu suất cao và cộng đồng hỗ trợ lớn, là lựa chọn lý tưởng cho các ứng dụng web cần lưu trữ dữ liệu có cấu trúc như thông tin người dùng, khóa học, bài thi, kết quả, và các dữ liệu liên quan khác.

Để tương tác với MySQL, dự án sử dụng \textbf{Spring Data JPA} kết hợp với \textbf{Hibernate} làm nhà cung cấp JPA (Java Persistence API). Đây là một giải pháp \texttt{Object-Relational Mapping (ORM)}, cho phép ánh xạ các đối tượng Java (Java objects) thành các bảng trong cơ sở dữ liệu quan hệ. Thay vì phải viết các câu lệnh SQL thủ công, lập trình viên có thể thực hiện các thao tác \texttt{CRUD} (Create, Read, Update, Delete) thông qua các phương thức Java trực quan trên các đối tượng. Ví dụ, để lưu một đối tượng \texttt{User} mới, chỉ cần gọi phương thức \texttt{userRepository.save(user)}. ORM không chỉ giúp mã nguồn sạch sẽ, dễ bảo trì hơn mà còn cung cấp các tính năng nâng cao như quản lý giao dịch (transaction management), \texttt{lazy loading} và \texttt{caching} để tối ưu hóa hiệu suất.

\subsection{Quản lý file với Cloudinary: Lưu trữ media trên cloud}

Để quản lý các file media như hình ảnh và âm thanh phục vụ cho các câu hỏi TOEIC, hệ thống tích hợp \textbf{Cloudinary} - một dịch vụ quản lý media trên cloud. Thay vì lưu trữ file trên server local, giải pháp này mang lại nhiều ưu điểm:

\begin{itemize}
    \item \textbf{Hiệu suất cao:} File được phục vụ từ CDN toàn cầu, đảm bảo tốc độ tải nhanh từ mọi vị trí địa lý.
    \item \textbf{Dung lượng không giới hạn:} Không bị ràng buộc bởi dung lượng ổ đĩa server, có thể lưu trữ hàng nghìn file audio và hình ảnh.
    \item \textbf{Tự động tối ưu:} Cloudinary tự động tối ưu hóa định dạng và chất lượng file dựa trên thiết bị và trình duyệt của người dùng.
    \item \textbf{Bảo mật:} Hỗ trợ HTTPS và các cơ chế bảo mật tiên tiến, đảm bảo an toàn cho nội dung.
    \item \textbf{API đơn giản:} Cung cấp API RESTful dễ sử dụng cho việc upload, quản lý và xóa file.
\end{itemize}

Quy trình quản lý file trong hệ thống:
\begin{itemize}
    \item \textbf{Upload:} Khi admin tạo câu hỏi mới, file audio/hình ảnh được upload trực tiếp lên Cloudinary thông qua API.
    \item \textbf{Lưu trữ URL:} Hệ thống chỉ lưu URL trả về từ Cloudinary vào database thay vì file thực tế.
    \item \textbf{Hiển thị:} Frontend sử dụng URL này để hiển thị file trực tiếp từ Cloudinary CDN.
    \item \textbf{Quản lý:} Hệ thống có thể xóa file không còn sử dụng thông qua Cloudinary API.
\end{itemize}

\subsection{Quản lý phiên làm việc và phân quyền}

Ngoài việc xác thực, hệ thống còn phải thực hiện phân quyền chi tiết. Thông tin về vai trò của người dùng (ví dụ: \texttt{ROLE\_USER}, \texttt{ROLE\_ADMIN}) được mã hóa bên trong payload của JWT. Khi một yêu cầu được gửi đến, \textbf{Spring Security} sẽ can thiệp, giải mã token và kiểm tra xem vai trò của người dùng có đủ quyền để truy cập vào endpoint được yêu cầu hay không. Việc phân quyền này được cấu hình ở cấp độ API, đảm bảo rằng một người dùng thông thường không thể thực hiện các hành động của quản trị viên.

\section{Giao tiếp giữa Frontend và Backend}

Frontend (React SPA) và backend (Spring Boot) giao tiếp với nhau hoàn toàn thông qua các lời gọi API RESTful qua giao thức HTTP/HTTPS. Khi người dùng thực hiện một hành động trên giao diện, chẳng hạn như bấm nút "Nộp bài", ứng dụng React sẽ sử dụng một thư viện như \texttt{Axios} hoặc \texttt{fetch} API của trình duyệt để tạo một HTTP request (ví dụ: \texttt{POST /api/exams/submit}) đến backend. Dữ liệu gửi đi và nhận về thường ở định dạng JSON (JavaScript Object Notation), một định dạng trao đổi dữ liệu nhẹ và dễ xử lý bởi cả JavaScript phía client và Java phía server.

Việc sử dụng HTTPS là bắt buộc để mã hóa toàn bộ dữ liệu truyền đi giữa client và server, đảm bảo rằng các thông tin nhạy cảm như mật khẩu hay token không bị kẻ xấu nghe lén trên đường truyền. Kiến trúc tách biệt này cũng cho phép caching hiệu quả các yêu cầu \texttt{GET}, giảm tải cho server và tăng tốc độ phản hồi cho người dùng.

\section{Tài liệu API với Swagger/OpenAPI: Tăng cường hiệu quả cộng tác}

Để đảm bảo sự phối hợp nhịp nhàng giữa đội ngũ frontend và backend, dự án tích hợp \textbf{Swagger} (theo đặc tả OpenAPI) để tự động hóa việc tạo tài liệu API. Bằng cách thêm các chú thích (annotations) vào các controller và model trong mã nguồn Spring Boot, Swagger có thể tự động tạo ra một tài liệu API chi tiết, chính xác và luôn được cập nhật.

Swagger mang lại những lợi ích thiết thực:
\begin{itemize}
    \item \textbf{Tài liệu sống:} Tài liệu API được sinh ra trực tiếp từ mã nguồn, đảm bảo nó luôn phản ánh đúng trạng thái hiện tại của các API, loại bỏ vấn đề tài liệu lỗi thời.
    \item \textbf{Giao diện tương tác (Swagger UI):} Swagger cung cấp một giao diện web trực quan, nơi các lập trình viên frontend và tester có thể xem danh sách tất cả các endpoint, thông tin chi tiết về tham số, định dạng dữ liệu trả về và các mã lỗi. Quan trọng hơn, họ có thể thực thi các lời gọi API trực tiếp từ trình duyệt để thử nghiệm mà không cần viết thêm bất kỳ mã nào.
    \item \textbf{Thúc đẩy phát triển song song:} Ngay khi backend định nghĩa xong các endpoint, họ có thể chia sẻ tài liệu Swagger. Đội ngũ frontend có thể dựa vào đó để phát triển giao diện mà không cần phải chờ backend hoàn thành toàn bộ logic, giúp rút ngắn đáng kể thời gian phát triển chung.
\end{itemize}

\subsection{Tích hợp OpenAI API: Chatbot hỗ trợ học tập thông minh}

Dự án tích hợp \textbf{OpenAI API} để phát triển một chatbot hỗ trợ người học trong quá trình ôn luyện TOEIC. Chatbot sử dụng mô hình ngôn ngữ tiên tiến từ OpenAI để tương tác tự nhiên, hỗ trợ người dùng với các chức năng như:

\begin{itemize}
    \item \textbf{Giải thích đáp án:} Sau mỗi câu hỏi trong đề thi thử, người học có thể yêu cầu chatbot giải thích vì sao một đáp án đúng/sai.
    \item \textbf{Gợi ý mẹo làm bài:} Chatbot có thể đưa ra các mẹo làm bài hoặc chiến lược phân tích câu hỏi để tối ưu hóa thời gian làm bài thi TOEIC.
    \item \textbf{Trả lời câu hỏi về từ vựng và ngữ pháp:} Người học có thể hỏi nghĩa của từ, cách dùng ngữ pháp, hoặc xin ví dụ trong ngữ cảnh thực tế.
    \item \textbf{Giao tiếp tự nhiên:} Chatbot hỗ trợ trò chuyện giống người thật, tạo cảm giác học tập thân thiện và dễ tiếp cận hơn.
\end{itemize}

Việc tích hợp OpenAI API giúp tăng trải nghiệm người dùng, cá nhân hóa quá trình học, đồng thời thể hiện khả năng mở rộng và ứng dụng AI trong giáo dục của hệ thống.

\section{Tổng kết}

Kiến trúc hệ thống được lựa chọn là sự kết hợp hài hòa giữa các công nghệ hàng đầu và đã được kiểm chứng trong ngành công nghiệp phần mềm. Việc sử dụng \textbf{React} và \textbf{TypeScript} ở frontend tạo ra một giao diện người dùng nhanh, linh hoạt và dễ bảo trì. \textbf{Spring Boot} cùng kiến trúc \textbf{RESTful API} ở backend cung cấp một nền tảng dịch vụ vững chắc, an toàn và có khả năng mở rộng cao. Cơ chế xác thực bằng \textbf{JWT}, cơ sở dữ liệu \textbf{MySQL}, dịch vụ quản lý file \textbf{Cloudinary} và việc quản lý API bằng \textbf{Swagger} đều là những mảnh ghép quan trọng, góp phần tạo nên một hệ thống hoàn chỉnh và chuyên nghiệp.

Sự lựa chọn này không chỉ đáp ứng các yêu cầu chức năng hiện tại của ứng dụng học TOEIC trực tuyến mà còn đặt nền móng vững chắc cho sự phát triển và mở rộng trong tương lai. Kiến trúc tách biệt rõ ràng, tài liệu hóa tốt và việc áp dụng các tiêu chuẩn ngành sẽ giúp quá trình bảo trì, nâng cấp và tích hợp các tính năng mới trở nên dễ dàng và hiệu quả hơn, đảm bảo chất lượng và sự bền vững của sản phẩm.

\end{document}